\scanpage{545}
% E545-01
and

% E545-02
% E545-03
\noindent(57)\quad
\diagram{\begin{tikzcd}[ampersand replacement=\&,column sep=large]
\deleted{\unsure{S_0}}M_{\rho,\sigma}[0]
  \arrow[r,"\delta_{SM}"]
\& \gamma_{\rho,\sigma} \arrow[r,"\chi_\gamma"]
\& \widehat{\mathbb Z}^{*}
\end{tikzcd}}
i.e. $\chi_{SM}$, (or \unsure{simply|simply} $\chi$)

% E545-04
consider

% E545-05
% E545-06
\noindent(58)\quad
\diagram{\begin{tikzcd}[ampersand replacement=\&,column sep=huge]
\deleted{\unsure{S_0}}M_{\rho,\sigma}[0]
 \arrow[r,"{\varepsilon_{\rho,\sigma}=(\unsure{\varepsilon_\sigma},\varepsilon_\rho)}"]
\& \{\pm1\}\times\{\pm1\}
\end{tikzcd}}

% E545-07
and

% E545-08
% E545-09
\noindent(59)\quad
\diagram{\begin{tikzcd}[ampersand replacement=\&,column sep=huge]
\deleted{\unsure{S_0}}M_{\rho,\sigma}[0]
 \arrow[r,"{\varepsilon_{2,3}:x\longmapsto(\unsure{\varepsilon_2}(\chi(x)),\varepsilon_3(\chi(x)))}"]
\& \{\pm1\}\times\{\pm1\}
\end{tikzcd}}

% E545-10
and consider the subgroup of coincidences of the two preceding homomorphisms

% E545-11
% E545-12
% E545-13
% E545-14
\[
\begin{aligned}
(60)\quad
\deleted{\unsure{S_0}}M_{\rho,\sigma}^{\unsure{\mathrm{eff}}}[0]
 &\overset{\mathrm{dfn}}=
 \{x\in\unsure{S_0\{M_{\rho,\sigma}\}}\mid
 \varepsilon_{\rho,\sigma}(x)=\varepsilon_{2,3}(x)\}
\\
 &=\left\{\left(
 \underset{\substack{\rotatebox{90}{$\in$}\\ B_{\rho,\sigma}}}{u},
 \underset{\substack{\rotatebox{90}{$\in$}\\ N_\rho}}{a},
 \underset{\substack{\rotatebox{90}{$\in$}\\ Z_\sigma}}{b}
 \right)\ \middle|\
 \begin{gathered}
 \delta_B(u)=\delta_\rho(a)=\delta_\sigma(b)
 \quad\text{i.e. }\Delta\in\gamma_{\rho,\sigma}\\
 \varepsilon_\rho(a)=
 \deleted{\unsure{\varepsilon_3(\chi_B(u))}}
 \inserted{\varepsilon_3(\chi_\gamma(\Delta))},
 \quad\varepsilon_\sigma(b)=\varepsilon_2(\chi_\gamma(\Delta))
 \end{gathered}
 \right.
\end{aligned}
\]

% E545-15
We have a natural map

% E545-16
\noindent(61)\quad
\diagram{\begin{tikzcd}[ampersand replacement=\&,column sep=huge]
M_{\rho,\sigma}^{\mathrm{eff}}[0] \arrow[r]
 \& \mathcal G'\times\mathcal G\times\widehat{\mathbb Z}^{*}\\
(u,a,b)\arrow[r,mapsto] \& (\alpha,\beta,\mu)
\end{tikzcd}}
, with
$\begin{gathered}
\alpha=ua^{-1}\tau'^{\nu_\rho(a)}\\
\beta=ub^{-1}\\
\mu=\chi_B(u)=\chi_\rho(a)=\chi_\sigma(b)
\end{gathered}$

% E545-17
\margin{\uline{N.B.} we have identified $a\in N_\rho$, $b\in N_\sigma$
\unsure{with their images} $i_\rho(a)$, $i_\sigma(b)$ in $G_{\rho,\sigma}$.}

% E545-18
% E545-19
\margin{\uline{N.B.} We have set \deleted{\unsure{(61)}}
$\nu_\rho(a)\in\mathbb Z/2$ such that
$(-1)^{\nu_\rho(a)}=\varepsilon_\rho(a)$}

% E545-20
and we have done what was needed for it to be \unsure{injective}
and have as image $S\unsure{\mathfrak{g}|g}_{\rho,\sigma}^{\mathrm{eff}}$, hence induce a bijection

% E545-21
\noindent(62)\quad
\diagram{\begin{tikzcd}[ampersand replacement=\&,column sep=large]
M_{\rho,\sigma}^{\mathrm{eff}}[0] \arrow[r,"\sim"]
 \& S\unsure{\mathfrak{g}|g}_{\rho,\sigma}^{\mathrm{eff}}.
\end{tikzcd}}

% E545-22
% E545-23
\uline{Remark}. In fact \inserted{the formulas for} (61)
\unsure{define} even a map

% E545-24
\noindent(62)\quad
\diagram{\begin{tikzcd}[ampersand replacement=\&,column sep=large]
M_{\rho,\sigma}[0]\arrow[r]
\&\mathcal G'\times\mathcal G\times\widehat{\mathbb Z}^{*}
\end{tikzcd}}

% E545-25
and for a $(u,a,b)\in M_{\rho,\sigma}[0]$, its image $(\alpha,\beta,\mu)$
is related to $u\in B_{\rho,\sigma}$ by the formulas

% E545-26
\margin{544}

% E545-F01
\margin{\begin{center}
544\\
\texttt{https://grothendieck.umontpellier.fr}\\
Université de Montpellier, UMR CNRS 5149, Case courrier 051, Place Eugène Bataillon\\
34095 Montpellier cedex 5 - France
\end{center}}

\scanpage{546}
\mspage{576}

% E546-01
% E546-02
% E546-03
\[
(63)\quad
\left\{\begin{aligned}
u(\rho)&=\operatorname{int}(\alpha)(\rho)
&&\bigl(\inserted{\mathcal G-}\text{conjugate to }\rho^{\unsure{\varepsilon'}},
\text{ where }\varepsilon'=\varepsilon_\rho(a)\bigr)\\
u(\sigma)&=\varepsilon\operatorname{int}(\beta)(\sigma)
&&\bigl(\mathcal G\text{-conjugate to }\sigma^\varepsilon,
\text{ where }\varepsilon=\varepsilon_\sigma(b)\bigr)\\
&\deleted{\unsure{\alpha}}
\end{aligned}\right.
\]

% E546-04
(but without necessarily having
$\varepsilon'=\varepsilon_3(\mu)$, $\varepsilon=\varepsilon_2(\mu)$).

% E546-05
% E546-06
This suggests a definition of a group a little bigger than $B_{\rho,\sigma}$,
in which $B_{\rho,\sigma}$ \unsure{will be} \inserted{included}
as a \unsure{normal|$\triangleleft$} subgroup \unsure{of index|of index}
$1,2$ or $4$ --- and which will have the advantage, in the ``good cases'', of being the group
of adelic points of a \unsure{suitable} group scheme over $\mathbb Z$.

% E546-07
% E546-08
On the other hand, to \deleted{\unsure{kill}} eliminate the sign
$\varepsilon$ in (63), (which is present \unsure{even} if
$(u,a,b)\in M_{\rho,\sigma}[0]^{\mathrm{eff}}$), I feel like replacing $\beta$ by

% E546-09
\deleted{$\unsure{\beta'}=\tau'^{\nu}\beta'$, where
$\nu\in\mathbb Z/2$ \unsure{is related}}

% E546-10
\[
(64)\quad \beta'=\tau'^{\nu}\beta
\qquad\left(\text{where }\nu\in\mathbb Z/2\mathbb Z\text{ related by }
(-1)^\nu=\varepsilon\right)
\]

% E546-11
So that we will have

% E546-12
\[
(65)\quad \deleted{\unsure{\varepsilon|\chi}}u(\sigma)=\operatorname{int}(\beta')(\sigma)
\quad\text{i.e.}\quad [u,\sigma]=\beta'\sigma(\beta')^{-1}.
\]

% E546-13
This is related to the convention \unsure{that a|which a}
$u\in\Gamma_{\mathbb Q}^{\sim}$, which we were specifying by the pair
$(\alpha,\beta)\in\widehat{\Pi}_{0,3}'\times\widehat{\Pi}_{0,3}$,
will henceforth be specified by the pair
$(\alpha,\beta')\in\widehat{\Pi}_{0,3}\times\widehat{\Pi}_{0,3}'$, with

% E546-14
\[
(66)\quad\unsure{\chi}(\alpha)=\varepsilon_3(\mu),
\quad\chi(\beta')=\varepsilon_2(\mu)
\qquad\bigl(\text{where }\mu=\chi(u)\text{ is the multiplier}\bigr)
\]

% E546-15
% E546-16
% E546-17
which \deleted{gives} for the operation of $u$ on
\deleted{$\unsure{\mathcal G_{0,3}^{+\wedge}}$}

% E546-18
\[
\mathcal G_{0,3}^{+\wedge}
=\mathrm{SL}(2,\mathbb Z)^{\wedge\unsure{\textcircled{1}}}
=\mathrm{SL}(2,\mathbb Z)^{\wedge}/\pm1
\]

% E546-19
does not in fact change the particular formulas:

% E546-20
% E546-21
\[
(67)\quad\left\{\begin{aligned}
u(\rho)&=\operatorname{int}(\alpha).\rho,\quad
u(\sigma)=\operatorname{int}(\beta)\sigma
\quad(=\operatorname{int}(\beta').\sigma)\\
\alpha\rho(\alpha)^{-1}
&=\beta\underset{\deleted{\unsure{\varepsilon}}}{\unsure{\sigma}}(\beta)^{-1}\ell_1^\nu
\quad(=\beta'\sigma(\beta')^{-1}\ell_1^\nu)
\end{aligned}\right.
\qquad\nu=\frac{\mu-1}{2}
\]

% E546-22
\margin{545}

% E546-F01
\margin{\begin{center}
545\\
\texttt{https://grothendieck.umontpellier.fr}\\
Université de Montpellier, UMR CNRS 5149, Case courrier 051, Place Eugène Bataillon\\
34095 Montpellier cedex 5 - France
\end{center}}

\scanpage{547}

% E547-01
% E547-02
but for the operation in
\deleted{$\unsure{\mathcal G=\mathcal T_{1,1}}$}

% E547-03
$\mathcal T_{1,1}^{+\wedge}\simeq\mathrm{SL}(2,\mathbb Z)^{\wedge}$
itself gives simpler formulas, namely, instead of

% E547-04
\[
(68)\quad\left\{\begin{aligned}
u(\rho)&=\operatorname{int}(\alpha)\rho\qquad
u(\sigma)=\varepsilon.\operatorname{int}(\beta)\sigma\\
\alpha\rho(\alpha)^{-1}&=\varepsilon\beta\sigma(\beta)^{-1}\ell_1^\nu
\end{aligned}\right.,
\]

% E547-05
the nicer formulas

% E547-06
\[
(69)\quad\left\{\begin{aligned}
u(\rho)&=\operatorname{int}(\alpha)\rho,\qquad
u(\sigma)=\operatorname{int}(\beta')(\sigma)\\
\alpha\rho(\alpha)^{-1}&=\beta'\sigma(\beta')^{-1}\ell_1^\nu
\end{aligned}\right.,
\]

% E547-07
where the sign $\varepsilon=\varepsilon_2(\mu)$ has disappeared. It appears in the formulas

% E547-08
\[
(70)\quad\chi(\alpha)=\varepsilon_3(\mu),\quad
\chi(\beta')=\varepsilon_2(\mu).
\]

% E547-09
It seems to me I have enough to be able to take up again ``freehand'' the
preceding constructions and \unsure{complete} them with more adequate notations.

% E547-10
% E547-11
% E547-12
The group $\mathcal G'$
\inserted{quotient of $\mathcal T_{1,1}^{\wedge}\simeq\mathrm{GL}(2,\mathbb Z)^{\wedge}$}
is moreover taking on increasing importance in our calculations, it is the one that deserves
\deleted{\unsure{also}} the simple notation $\mathcal G$, the closed subgroup
generated by $\rho,\sigma$ becoming $\mathcal G^{+}$ --- it corresponds to
$\mathcal T_{1,1}^{+\wedge}\simeq\mathrm{SL}(2,\unsure{\widehat{\mathbb Z}|\mathbb Z})^{\wedge}$.

% E547-13
We will simply assume that $\tau\notin\mathcal G^{+}$
($\tau$ being the image of $\tau_\infty$), or what amounts to the
same, that

% E547-14
\[
(71)\quad\tau'\notin\mathcal G^{+}
\]

% E547-15
\margin{546}

% E547-F01
\margin{\begin{center}
546\\
\texttt{https://grothendieck.umontpellier.fr}\\
Université de Montpellier, UMR CNRS 5149, Case courrier 051, Place Eugène Bataillon\\
34095 Montpellier cedex 5 - France
\end{center}}

\scanpage{548}
\mspage{577}

% E548-01
% E548-02
It is ultimately this element, which normalizes both $L_\rho$ and $L_\sigma$
\deleted{\unsure{by}} \unsure{via} the remarkable formulas

% E548-03
\[
(*)\quad \tau'(\rho)=\rho^{-1},\qquad \tau'(\sigma)=\sigma^{-1}
\]

% E548-04
that takes precedence over $\tau$, which satisfies

% E548-05
\[
(**)\quad \tau(\varepsilon_0)=\varepsilon_0^{-1},\qquad
\tau(\varepsilon_1)=\varepsilon_1^{-1}
\quad\text{and even }\tau(\sigma)=\sigma^{-1}
\]

% E548-06
(which is obviously very \unsure{virtuous} !), \unsure{but} whose action on $\rho$

% E548-07
\[
\tau(\rho)=\sigma(\rho^{-1})
\]

% E548-08
\unsure{is} less convenient.

% E548-09
\margin{Nevertheless $\tau'$ has over $\tau$ the serious disadvantage of not being
an \unsure{element} of the group
$\widehat{\Pi}'_{0,3}=\{1,\tau\}.\widehat{\Pi}_{0,3}$, where we have the rigidity relation
$\widehat{\Pi}_{0,3}^{\prime\sigma}=\widehat{\Pi}_{0,3}^{\unsure{\prime}\rho}=\{1\}$
--- whereas in $\{\unsure{1},\tau'\}.\widehat{\Pi}_{0,3}$, the centralizer of
$\sigma$ is not reduced to $1$, but is equal to $\{1,\tau'\}$
Moreover, $\tau'$ has the disadvantage of not being in
$\unsure{\mathcal M}_{0,3}^{\unsure{!|\prime}}$ \dots}

% E548-10
\[\overline{\underline{\mathrm V}}\]

% E548-11
\margin{\uline{Preliminary}\par\uline{data}}

% E548-12
$1^\circ)$ So let $\mathcal G$ be a profinite group, equipped with a surjective homomorphism

% E548-13
\[\unsure{=|\to}\]

% E548-14
\noindent(1)\quad
\diagram{\begin{tikzcd}[ampersand replacement=\&,column sep=large]
\mathcal T_{1,1}^{\unsure{*|+}\wedge}\simeq\mathrm{GL}(2,\mathbb Z)^{\wedge}
\arrow[r,"\text{epimorphism}"] \& \mathcal G
\end{tikzcd}},

% E548-15
which we will express by giving the three topological generators
$\rho,\sigma,\tau'\in\mathcal G$, satisfying

% E548-16
\[
(2)\quad \rho^6=\sigma^4=\tau'^2=1,\quad
\begin{gathered}\rho^3=\sigma^2\\\text{(denoted }-1\text{)}\end{gathered},\quad
\tau'\rho\tau'^{-1}=\rho^{-1},\quad
\tau'\sigma\tau'^{-1}=\sigma^{-1}.
\]

% E548-17
We denote by $\mathcal G^+$ the closed subgroup generated by $\rho,\sigma$,
which is thus the image of

% E548-18
\[
\mathcal T_{1,1}^{+\wedge}\simeq\mathrm{SL}(2,\mathbb Z)^{\wedge},
\]

% E548-19
\margin{547}

% E548-F01
\margin{\begin{center}
547\\
\texttt{https://grothendieck.umontpellier.fr}\\
Université de Montpellier, UMR CNRS 5149, Case courrier 051, Place Eugène Bataillon\\
34095 Montpellier cedex 5 - France
\end{center}}